%% Écriture d'un torseur d'action mécanique : résultante (X, Y, Z) à gauche,
%% moment (L, M, N) à droite, point de réduction A écrit sous l'accolade fermante.
\documentclass[tikz,border=4pt]{standalone}
\usepackage{liaisons}
\usetikzlibrary{positioning}
\begin{document}
\begin{tikzpicture}[x=1cm,y=1cm,
  bulle/.style={draw=black!45, fill=black!3, rounded corners=4pt, inner sep=4pt,
                font=\small, align=center},
  lien/.style={-{Stealth[length=5pt]}, line width=0.6pt, black!60}]

% --- macro locale : accolades à la hauteur des trois lignes ------------
\newcommand\gabarit{\vphantom{\begin{matrix} X_{1\rightarrow 2}\\[2pt] Y_{1\rightarrow 2}\\[2pt] Z_{1\rightarrow 2}\end{matrix}}}

% --- le torseur : deux colonnes entre accolades ------------------------
\node[font=\small] (LC) at (3.5,0)
     {$\begin{matrix} X_{1\rightarrow 2}\\[2pt] Y_{1\rightarrow 2}\\[2pt] Z_{1\rightarrow 2}\end{matrix}$};
\node[font=\small, right=0.55cm of LC] (RC)
     {$\begin{matrix} L_{A\,1\rightarrow 2}\\[2pt] M_{A\,1\rightarrow 2}\\[2pt] N_{A\,1\rightarrow 2}\end{matrix}$};
\node[font=\small, left=2pt of LC]  (BL) {$\left\{\gabarit\right.$};
\node[font=\small, right=2pt of RC] (BR) {$\left.\gabarit\right\}$};
\node[font=\small, left=2pt of BL] {$\{\tau_{1\rightarrow 2}\} =$};
\node[font=\small, anchor=north west, inner sep=1pt] (A) at ($(BR.south east)+(-2pt,4pt)$) {$A$};

% --- ce que contient chaque colonne ------------------------------------
\node[bulle, anchor=south] (RES) at (2.45,1.15) {\textbf{résultante} (la force), en N};
\node[bulle, anchor=south] (MOM) at (6.45,1.15) {\textbf{moment}, en N$\cdot$m};
\draw[lien] (RES.south) -- (LC.north);
\draw[lien] (MOM.south) -- (RC.north);

% --- le point de réduction ---------------------------------------------
\node[bulle, anchor=north] (PR) at (3.6,-1.25)
     {\textbf{point de réduction} : le point où est exprimé le moment};
\draw[lien] (PR.north east) -- ($(A.south)+(0,-0.04)$);
\end{tikzpicture}
\end{document}
